% IEEE conference paper template (IEEEtran, ships with TeX Live).
% Compile: latexmk -pdf ieee-paper.tex   (runs bibtex automatically)
% Needs refs.bib in the same directory.
\documentclass[conference]{IEEEtran}

\usepackage{amsmath,amssymb}
\usepackage{graphicx}
\usepackage{booktabs}
\usepackage{cite}

\begin{document}

\title{Paper Title Goes Here:\\ A Subtitle If Needed}

\author{
  \IEEEauthorblockN{First Author}
  \IEEEauthorblockA{Department of Electrical Engineering\\
    University Name\\
    City, Country\\
    first.author@example.edu}
  \and
  \IEEEauthorblockN{Second Author}
  \IEEEauthorblockA{Research Lab\\
    Company Name\\
    City, Country\\
    second.author@example.com}
}

\maketitle

\begin{abstract}
One paragraph, roughly 150--250 words. State the problem, the approach,
and the headline result with a number. Avoid citations and undefined
acronyms in the abstract.
\end{abstract}

\begin{IEEEkeywords}
keyword one, keyword two, keyword three
\end{IEEEkeywords}

\section{Introduction}
Motivate the problem and summarize contributions. Cite prior work
like this~\cite{shannon1948communication} or grouped~\cite{he2016resnet,vaswani2017attention}.

The contributions of this paper are:
\begin{itemize}
  \item First contribution, stated as a claim.
  \item Second contribution.
  \item Third contribution, with the key metric.
\end{itemize}

\section{Related Work}
Group prior work by theme, not chronologically. Deep residual
networks~\cite{he2016resnet} and attention
mechanisms~\cite{vaswani2017attention} are typical anchors.

\section{Method}
\subsection{Problem Formulation}
Inline math: let $x \in \mathbb{R}^n$ denote the input. Display math
with a number:
\begin{equation}
  y = \sigma\!\left( W x + b \right),
  \label{eq:model}
\end{equation}
and refer back to it as \eqref{eq:model}. For multi-line derivations:
\begin{align}
  P_\mathrm{out} &= P_\mathrm{in} \, G_1 G_2 \label{eq:gain}\\
                 &= P_\mathrm{in} \cdot 10^{(G_{1,\mathrm{dB}} + G_{2,\mathrm{dB}})/10}.
\end{align}

\subsection{Figure Example}
Reference figures as Fig.~\ref{fig:arch}. Replace the placeholder
\verb|\rule| with \verb|\includegraphics[width=\columnwidth]{arch.pdf}|.
\begin{figure}[t]
  \centering
  \rule{0.8\columnwidth}{3cm}% placeholder box
  \caption{System architecture. Replace this placeholder with an
  \texttt{includegraphics} call.}
  \label{fig:arch}
\end{figure}

\section{Results}
Reference tables as Table~\ref{tab:results}.
\begin{table}[t]
  \caption{Comparison With Prior Work}
  \label{tab:results}
  \centering
  \begin{tabular}{lcc}
    \toprule
    Method & Accuracy (\%) & Params (M) \\
    \midrule
    Baseline~\cite{he2016resnet} & 92.1 & 25.6 \\
    Prior art~\cite{vaswani2017attention} & 93.4 & 65.0 \\
    \textbf{Ours} & \textbf{94.7} & 24.1 \\
    \bottomrule
  \end{tabular}
\end{table}

\section{Conclusion}
Restate the contribution and result; one sentence of future work.

\section*{Acknowledgment}
Optional. Funding sources and colleagues.

\bibliographystyle{IEEEtran}
\bibliography{refs}

\end{document}
